\chapter{Tabla consolidada de causas raíz} \label{ape:causas-raiz}

Este apéndice recoge, de forma consolidada, los problemas documentados a lo largo del trabajo con su causa raíz y su categoría. Sirve como referencia rápida para reproducir el razonamiento de los capítulos~\ref{cap3} y~\ref{cap4}.

\section{Capa host: instalación del driver 340.108}

\begin{table}[htbp]
\centering
\footnotesize
\setlength{\tabcolsep}{3pt}
\begin{tabular}{|c|L{3.5cm}|L{5.5cm}|L{3.8cm}|}
\hline
\# & \textbf{Problema} & \textbf{Causa raíz} & \textbf{Categoría} \\
\hline
1 & Conflicto de configuración & Residuos de configuración VFIO reservando las GPUs & Configuración del host \\
\hline
2 & Cabeceras (\textit{headers}) no encontradas & Cambio de nomenclatura de paquetes en Proxmox VE 9 (\texttt{pve-*} $\rightarrow$ \texttt{proxmox-*}) & Empaquetado / doc. desactualizada \\
\hline
3 & Kernel demasiado nuevo & $\sim$7 años de diferencia entre el driver (2019) y el kernel por defecto (2025--2026) & Obsolescencia driver \textit{legacy} \\
\hline
4 & Desajuste de versión de gcc & Diferencia de \textit{patch version} entre el gcc del kernel y el del sistema & Comprobaciones estrictas del instalador NVIDIA \\
\hline
5--6 & Cabeceras del compilador no resueltas & Interacción entre \texttt{-nostdinc} y múltiples versiones de gcc coexistentes & Configuración del \textit{toolchain} \\
\hline
7 & Tipos del driver no definidos & Fallo del mecanismo \texttt{conftest} al sobrescribir flags manualmente & Metodología de compilación \\
\hline
8--9 & Necesidad de parches de terceros & API del kernel completamente distinta a la esperada por el código de 2019 & Sin mantenimiento oficial \\
\hline
10 & Opción de compilación no reconocida & Parche comunitario orientado a gcc 15+ aplicado sobre gcc 12.4 & Desfase entre parche y \textit{toolchain} \\
\hline
\end{tabular}
\caption{Causas raíz de la instalación del driver en el host}
\label{tab:causas-host}
\end{table}

\section{Capa contenedor/CUDA}

\begin{table}[htbp]
\centering
\footnotesize
\setlength{\tabcolsep}{3pt}
\begin{tabular}{|c|L{3.5cm}|L{5.5cm}|L{3.8cm}|}
\hline
\# & \textbf{Problema} & \textbf{Causa raíz} & \textbf{Categoría} \\
\hline
1 & Plantilla no encontrada & La nomenclatura de la versión de \textit{build} cambia con el tiempo en el índice de Proxmox & Doc. desactualizada \\
\hline
2 & Compilador no soportado (instalador CUDA) & CUDA 6.5 (2014) declara soporte hasta gcc 4.8; el sistema tiene gcc 12.2 & Obsolescencia \textit{legacy} \\
\hline
3 & Instalación silenciosa fallida sin aviso claro & El meta-instalador \texttt{--silent} oculta el fallo real de un subinstalador anidado & Instalador poco transparente \\
\hline
4 & \texttt{InstallUtils.pm} no localizado & Ruta \texttt{@INC} rota en el flujo de autoextracción \textit{makeself} & Fragilidad del empaquetado \\
\hline
5 & El script auxiliar de samples falla & Asume una estructura (\texttt{../samples/}) distinta de la del paquete (\texttt{./cuda-samples/}) & Inconsistencia del paquete \\
\hline
6 & Front-end de \texttt{nvcc} incompatible con los \textit{headers} modernos & Analizador C++ de \texttt{nvcc} 6.5 (previo a C++11) frente a la glibc de Debian 12 & Obsolescencia \textit{toolchain} \\
\hline
\end{tabular}
\caption{Causas raíz de la capa contenedor/CUDA}
\label{tab:causas-contenedor}
\end{table}

